\title{Improving alignment of dialogue agents via targeted human judgements}

\begin{abstract}
We present Sparrow, an information-seeking dialogue agent trained to be more helpful, correct, and harmless compared to prompted language model baselines.  We use reinforcement learning from human feedback to train our models with two new additions to help human raters judge agent behaviour.  First, to make our agent more helpful and harmless, we break down the requirements for good dialogue into natural language rules the agent should follow, and ask raters about each rule separately.  We demonstrate that this breakdown enables us to collect more targeted human judgements of agent behaviour and allows for more efficient rule-conditional reward models.  Second, our agent provides evidence from sources supporting factual claims when collecting preference judgements over model statements.  For factual questions, evidence provided by Sparrow supports the sampled response 78\% of the time.  Sparrow is preferred more often than baselines while being more resilient to adversarial probing by humans, violating our rules only 8\% of the time when probed.  Finally, we conduct extensive analyses showing that though our model learns to follow our rules it can exhibit distributional biases.
\end{abstract}

\section{Introduction}
\label{sec:intro}

Many deployed machine learning systems operate in settings in which there is no program that computes the system's objective. This is true not only of many natural language tasks, but also of robotics and other tasks where only some aspects of safe behaviour and task completion can be specified a priori. This lack of programmatic reward motivates reinforcement learning from human feedback (RLHF) where human judgements of behaviour are a critical component of the training process. However, human supervision works only if the humans are well-informed and motivated, and if the data collection setup is robust to human errors.

In this paper we study the use of human judgements as rewards for the task of helpful, correct, and harmless \emph{information-seeking dialogue}, defined as a conversation between a human user and a dialogue agent with the goal of providing answers to questions and follow-ups posed by the user~\citep{zamani2022conversational}. Dialogue allows users to naturally communicate their intentions to the agent. Dialogue is also very general, posing both opportunities for novel behaviours and many concrete harms that must be addressed \citep{bender2021dangers, weidinger2021ethical}. By focusing on information-seeking dialogue, the context and criteria for success are better-defined (e.g. \emph{Was the information provided?}) than for so-called \emph{chit-chat} dialogue, and better-defined contexts make it easier to define harms. We call the resulting model Sparrow.

Our primary contributions are:

\begin{enumerate}
\item \textbf{Targeted human judgements of specific rules}: We guide and elicit targeted judgements from human annotators by asking about violations for a number of rules such as "Do not make statements which are threatening" or "Do not offer financial advice" (see \ref{tab:harm-rules-short}). This lets us characterise failures of the model, train targeted classifiers, and guide humans towards probing failure modes of interest. This extends previous probing methods that focus on simply safe/unsafe labels~\citep{xu2021bot} or broad notions of harm~\citep{askell2021general, bai2022training}.

\item \textbf{Multi-objective RLHF to maximise preference rates and minimise rule violations}: We successfully combine a variety of techniques to train a single unified model. We show that by combining targeted rule judgements and preference judgements with RLHF, we can train a model that is preferred to baselines based on prompting, reranking or supervised learning alone (\ref{fig:pareto_plot}). Simultaneously, Sparrow is much more resilient to adversarial attacks by humans than our baselines, breaking the targeted rule in only 8\% of probe conversations.

\item \textbf{Inline evidence to improve correctness and verifiability}: We adapt and extend the methods of GopherCite \citep{menick2022teaching} to the interactive dialogue setting, while demonstrating performance similar to GopherCite on single-turn QA tasks.  When Sparrow provides answers with evidence, those answers are supported and plausible 78\% of the time, a significant improvement over our prompted baselines. Providing evidence helps raters verify claims.

\item \textbf{Detailed analyses of the resulting dialogue agent}:  In particular, we highlight our analysis of the impact of our methods on the \emph{distributional} properties of the resulting RL policy, as our mitigations address only \emph{instance harms}~\citep{weidinger2021ethical}.  Our findings show that our methods, although they improve rule following, can amplify distributional fairness concerns.
\end{enumerate}

Our work shares many features with other dialogue systems such as LaMDA~\citep{thoppilan2022lamda}, the Anthropic assistant~\citep{askell2021general, bai2022training}, and SeeKeR~\citep{shuster2022seeker}. LaMDA also collects annotations for individual rules, but does not use per-rule labels when mitigating or evaluating rule violations, and uses supervised learning and ranking rather than reinforcement learning. We borrow the \emph{helpful}, \emph{honest}, and \emph{harmless} (HHH) decomposition of \citet{askell2021general}, but use \emph{correct} instead of \emph{honest} for now as our methods do not address honesty directly.  \citet{bai2022training} uses reinforcement learning from human preferences to train a dialogue agent to be helpful and harmless, but does not break rules down further for humans, trains a single reward model to represent all human feedback, and does not incorporate external evidence. SeeKeR, LaMDA, and BlenderBot 3 use a similar knowledge retrieval mechanism where a generated search query is used to retrieve information on which the response is conditioned, but SeeKeR does not show the retrieved information to raters during evaluation, and none of these use RL.

Although the mechanisms introduced here are a useful starting point for robust alignment of models, we point out several areas of necessary future work. Besides its role as a task, we believe dialogue is a flexible medium through which various sources of evidence and instructions can be combined to help humans evaluate agent behaviour. In the future, this might include methods such as debate \citep{irving2018debate} where agents present arguments for and against their previous outputs to assist with human judgement.

\section{Methods} \label{sec:methods}
Starting with Dialogue Prompted Chinchilla 70B (DPC)~\citep{hoffmann2022training} described in \ref{sec:prompting}, we gather human data for rule violations and per-turn response preferences (\ref{sec:eval_start}). This data is used to train preference reward models (\emph{preference RMs}) and a rule reward model (\emph{rule RM}) that predicts whether a given rule was violated (\ref{sec:reward-models}). We use reinforcement learning with advantage actor-critic (A2C)~\citep{mnih2016a3c} to train, initialised from our DPC base model. We jointly optimise for the rule violation rate estimated by the rule RM and per-turn response preferences estimated by preference RMs (\ref{sec:rl}). We continuously expand our set of ratings through data collections with improved models, and in turn improve our models with more data (\ref{fig:training_pipeline}) following \citet{stiennon2020learning}. In addition to RL, we also employ our reward models for reranking at test-time (\ref{sec:rerank}) to further improve performance.

\subsection{Defining rules} \label{sec:rules}

Starting with our high-level goals of \emph{helpful}, \emph{correct}, and \emph{harmless} dialogue, we divide each goal into more detailed rules, shown in \autoref{tab:harm-rules-short}, for use in rule-based adversarial probing and rule-conditional classification. Helpfulness rules include answering user questions, staying on topic, and avoiding common problems such as repetition, and are combined with an overall \emph{per-turn response preference} in \ref{sec:eval_start}.
Correctness rules cover types of incorrect statements which raters might not otherwise penalise, such as the agent claiming to have a physical body or take real-world actions, and are combined with the evidence-specific rewards in \ref{sec:evidence}.
Both helpfulness and correctness rules are frequently violated by our baseline model.

Though prior work has demonstrated that language and dialogue models can output harmful language \citep{rae2021gopher,gpt3,dinan-etal-2019-build}, our baseline produced language we consider harmful only infrequently or under adversarial behaviour by users.  Consequently, instead of writing rules based on example failures, we consulted existing literature to identify potential failure modes, then wrote rules and sought examples where our model would fail to follow our rules.

We designed our rule set to test our methodology with a set of representative requirements for information-seeking agents; we did not aim for completeness in our rule set. In particular, we focused on harms which can be encoded in natural-language rules and mitigated using RL from human feedback, as other goals such as social, linguistic or environmental justice~\citep{bender2021dangers} require mitigation strategies outside the scope of this paper.  Broadly, we identified candidate risks which fall under discrimination, exclusion, toxicity, misinformation, and human-computer interaction harms in the taxonomy proposed by \citet{weidinger2021ethical}.  Prior work has argued that broad umbrella terms like ``toxicity'' can be ambiguous~\citep{welbl-etal-2021-challenges-detoxifying,vidgen-etal-2019-challenges,banko-etal-2020-unified} so we base our rules on more fine-grained definitions including the online harmful language taxonomy from \citet{banko-etal-2020-unified}, and definitions of microaggressions from \citet{breitfeller-etal-2019-finding}.  For rules which pertain to providing legal, financial, and medical advice, we consulted an in-house legal expert.  \ref{app:harm-rules} lists our rules and specific resources used to craft each rule.  While we put extensive thought into our initial rule set, we emphasise that they are not comprehensive and require substantial expansion and refinement before real-world usage.

Our rules resemble the safety objectives in \citet{thoppilan2022lamda}, but were crafted with our annotation process in mind.  In particular, within a single task, we ask annotators to consider a different rule for each dialogue they generate and annotate.  To help annotators comprehend different rules quickly, our rules are designed to be short and standalone (so that an annotator can understand an individual rule without any other context).

\begin{center}

\end{center}

\subsection{Generating dialogue turns}

\paragraph{Prompting for dialogue}\label{sec:prompting}
Following \citet{rae2021gopher}, we construct a dialogue agent by combining  Chinchilla-70B~\citep{hoffmann2022training} with a hand-authored prompt that demonstrates good behaviour in a dialogue between two participants: \texttt{User} and \texttt{Agent}. For a dialogue agent with evidence, we introduce two new participants: \texttt{Search Query}, which generates a search query; and \texttt{Search Result} which adds the evidence retrieved from Google Search based on the \texttt{Search Query} turn, similar to \citet{lazaridou2022internet}.
To generate \texttt{User}, \texttt{Search Query}, and \texttt{Agent} turns, the prompt, dialogue history, and participant name are concatenated and form the context for which completion is sampled using nucleus sampling \citep{holtzman2019nucleus}. \texttt{Search Result} turns are constructed by querying Google Search and scraping the returned search results, as described in \ref{itm:retrieval}. \ref{fig:evidence_text_to_ui} illustrates how the dialogue transcript is formatted into an LM context, and rendered when displayed to humans.

We iterated on the prompt, informed by behaviours seen during evaluation. Throughout the paper DPC (Dialogue-Prompted Chinchilla) refers to unmodified Chinchilla with our final prompt (\ref{app:prompt}); other models are prompted in the same way, unless indicated otherwise. \ref{fig:rl-ablations} characterises the effects of using a different prompt from \citet{rae2021gopher}.

\paragraph{Selecting whether to use evidence}\label{itm:choose-search}
Throughout this and the following sections we use the following nomenclature to refer to different methods for determining whether to use evidence:

\begin{itemize}
\item \emph{always search}: A model that is forced to produce a \texttt{Search Query} turn and condition on the \texttt{Search Result}.
\item \emph{never search}: A model that is forced to produce an \texttt{Agent} turn without evidence.
\item \emph{choose search}: The selection of whether to search or not is made by computing the log likelihood for the roles \texttt{Search Query} and \texttt{Agent} following the dialogue context. The role with the higher log likelihood is chosen to continue the dialogue, which determines whether we use evidence retrieved from Google Search in the response or not.
\item \emph{@N}: Instead of choosing whether to search or not, we produce \emph{N} responses: half the responses are produced by generating search queries and conditioning on \texttt{Search Results}, the other half are generated without evidence. Whether the final response uses evidence is determined by reranking with reward models, as described in \ref{sec:rerank}.

\end{itemize}

\subsection{Human data collection} \label{sec:eval_start}

Similar to \citet{ouyang2022instructgpt,stiennon2020learning} and others, our method involves a continuous cycle of evaluation and training as illustrated in \ref{fig:training_pipeline}. 
We start with DPC described in \ref{sec:prompting}, as the initial dialogue agent. We then ask human participants to interact with the agent in two main settings: \emph{per-turn response preference} and \emph{adversarial probing}.

\paragraph{Per-turn response preference} \phantomsection\label{itm:response-preference}
In this task, human raters are given an incomplete dialogue and multiple possible statements to continue the dialogue, each corresponding to a different sample or model. The human raters select the response that they think is best (\ref{fig:screenshot_preference_rating}).
In contrast to \citet{askell2021general}, a model generates both the \texttt{User} and \texttt{Agent} turns, and in both cases the human raters are asked to select the best response. The selected response is then used to continue the dialogue. Per-turn response preference data lets us estimate a \emph{preference rate} which measures how frequently a model is preferred over one or more competing models. When responses are combined with supporting evidence, human raters give additional per-response feedback, as described in \ref{sec:evidence}.

\paragraph{Adversarial probing} \phantomsection\label{itm:adversarial-probing}
In this task, we show participants one of the rules, and they are instructed to have a conversation that leads the model to break the rule. Following the conversation, the same participant then indicates whether the model followed the rule or not. Instructing participants to focus on specific rules rather than a general rule allows us to target and improve on specific failure modes (\ref{sec:general}).
Collecting many dialogues of this form let us estimate a rule violation rate under human adversarial probing. This approach extends ideas from \citet{xu-etal-2021-detoxifying} to fine-grained rules. Representative images of the per-turn response preference and adversarial probing tasks are included in \ref{app:ui-screenshots}.

\paragraph{Training and evaluation pipeline} Adversarial probing and per-turn response preference allow us to improve the model. Adversarial probing is used to assess how vulnerable the model is to exhibiting bad behavior and the response preference rate is used as a measure for helpfulness (see axes in \ref{fig:pareto_plot}).
From our rule violation data, we train a \emph{Rule RM} (reward model) that predicts human judgment of rule violation. The preference data is used to train Elo \emph{Preference RMs} as a proxy for helpfulness (\ref{sec:reward-models}).
We then use both the Rule RM and the Preference RMs to improve the agent via reranking (\ref{sec:rerank}) and RL (\ref{sec:rl}).

\paragraph{Data quality} Even after appropriate consideration, raters do not always agree about whether a rule was violated by Sparrow in a given conversation. Raters often lack the knowledge or context to determine whether statements are faithfully-grounded with evidence and some of the requirements of good behavior are ambiguous or under-specified. We ask our participants to complete an interactive click-through tutorial before the actual tasks to assist the raters with task understanding, and used comprehension checks to improve data quality (see \ref{app:human-data}). Despite the remaining disagreements inherent to human judgement, we believe that higher per-turn preference rates and lower rule violation rates correspond to improvements to our model.

\paragraph{Annotator well-being} The details of our study design, including compensation rates, were reviewed by our independent ethical review committee. All participants provided informed consent prior to completing tasks and were reimbursed for their time. It is our policy that researchers must pay workers/participants at least the living wage for their location. Because some of our rules refer to sensitive topics and could plausibly cause psychological or emotional harm to our annotators \citep{dang2018but,steiger2021psychological}, we monitored rater well-being through a well-being survey. We set data budgets for sensitive topics and structured rating tasks such that raters were allowed to skip tasks and rules for well-being reasons without penalty at any point. A summary of well-being survey results is available in \ref{app:annotator-wellbeing}, along with statistics capturing the broad demographics of raters that participated.

\paragraph{Related work} Our human data collection protocols share some commonalities with those used to train and evaluate LaMDA~\citep{thoppilan2022lamda}, the Anthropic assistant~\citep{askell2021general, bai2022training}, WebGPT \citep{nakano2021webgpt}, and BlenderBot 3~\citep{blenderbot3}. BlenderBot 3 collects non-adversarial open-domain short conversations, soliciting binary per-turn feedback and suggestions for an improved response. LaMDA collects dialogues in both adversarial and non-adversarial settings. The transcripts are labeled separately, and used for classifier training as well as evaluation against quality and safety metric. Neither BlenderBot 3 nor LaMDA collect preference ratings between model responses for training or evaluation, and opt instead for absolute score-based approaches. The Anthropic assistant uses a unified protocol in which user turns are human-generated and agent turns are chosen from two possible responses. Their data collection follows one of two modes: having raters either pick the best response, or the worst response at each turn --- these correspond in purpose to our user preference and adversarial collections, respectively. In common with WebGPT, a key component of our evaluation set-up is that Sparrow surfaces evidence (\ref{sec:evidence}) for its claims in the form of excerpts from the web; this allows the raters to more easily verify its claims without needing to do independent research.

\subsection{Evidence} \label{sec:evidence}

We train our model to search the internet in order to provide more correct responses. This mechanism also allows for temporal generalisation beyond a static parametric model \citep{liska2022streamingqa,lewis2020retrieval,borgeaud2022retro,shuster2022seeker}. In our user interface, we display the evidence used by the model next to the model's response to assist the rater in appraising whether the model's response is correct (\ref{fig:evidence_text_to_ui}).
Supporting model responses with evidence \citep{menick2022teaching} serves as a type of explanation  \citep{ras2022explainable}, providing an insight into the external information the model was provided when generating the answer. This allows raters to better assess factual accuracy and affords end-users greater trust in the model (\ref{sec:rater-perception-eval}).

\paragraph{Learning to search} To learn how to search and when to use the evidence, we train a preference model from human judgements on samples from existing models (DPC or earlier versions of Sparrow). We bootstrap from an initial evidence-supported dialogue model by prompting~\citep{lazaridou2022internet,menick2022teaching}. We incorporate evidence into the dialogue framework by introducing two participants into the dialogue prompt: \texttt{Search Query} and \texttt{Search Result}. \ref{app:prompt_evidence} details the prompt and baseline model.

Response preferences are collected over four-statement comparisons; two responses are sampled without evidence from agents with the non-evidence prompt (\ref{app:prompt}), while the other two agents first generate search queries, obtain search results, and condition on the evidence to produce their responses. The rater's choice between these four options provides signal both for the overall quality of the response and search query (if used), and for the decision to display evidence or not.

\paragraph{Retrieval} \phantomsection\label{itm:retrieval}

The \texttt{Search Result} turn is constructed by retrieving Google Search results for a \texttt{Search Query} sampled from Sparrow. We scrape the returned HTML web pages and truncate a fragment of up to 500-characters around the search engine-provided snippet for each result (\ref{app:retrieval}). A \texttt{Search Result} turn contains a single scraped fragment and is added to the dialogue context for the \texttt{Agent}. This turn is displayed to the raters as evidence quoted from the web (\ref{fig:evidence_text_to_ui}). 

\paragraph{Collecting human feedback} \phantomsection\label{itm:evidence-app}

Given a model that can optionally search, we aim to assess two properties. First, how often does the model provide evidence when making a factual claim? Second, how often does the evidence (when provided) support the claims of the model? To make these assessments, we ask raters additional questions about the dialogue when collecting response preferences. In particular, raters are asked the following questions:

Before seeing possible responses (see \ref{fig:screenshot_question_annotation}):

\begin{itemize}
\item{Should the AI search the internet to support its response?}
\end{itemize}

For each response with evidence, individually (see \ref{fig:screenshot_evidence_annotation}):

\begin{itemize}
\item{Is the response plausible (reasonable, on topic, could be true)?}
\item{Is the response supported by the provided evidence from the internet? (i.e.\ the evidence convinces you that the answer is correct)}
\end{itemize}

For each response without evidence, individually (see \ref{fig:screenshot_non_evidence_annotation}):

\begin{itemize}
\item{Is this response plausible (reasonable, on topic, could be true)?}
\item{Could this response be supported by quoting facts from the internet?}
\end{itemize}

Responses to these questions let us investigate how often the model provides evidence when needed, and how often it successfully makes claims that are supported by evidence. Measuring and optimising towards the supportedness of evidence is important for assessing and increasing the rate at which responses are faithfully-grounded in external knowledge, and reducing the problem of hallucinations~\citep{dziri2022hallucination}. We ask the above questions (see \ref{fig:screenshot_preference_rating}) for every response option as part of the response preference task, before the selection of the best option (see \ref{sec:eval_start}). 

\subsection{Reward models} \label{sec:reward-models}

We train two types of reward models separately, both fine-tuned from  Chinchilla 70B:

\begin{itemize}
    \item The \textbf{Response Preference Reward Model (Preference RM)} scores responses according to human preferences between candidate responses. 
    \item The \textbf{Rule Violation Reward Model (Rule RM)} estimates the probability that Sparrow breaks a rule in a given dialogue.
\end{itemize}

Response preference data (\ref{itm:response-preference}) allows us to train a Preference RM that for each response predicts an \emph{Elo} preference score such that the softmax over the scores predicts the preference probability, following~\citep{stiennon2020learning,menick2022teaching,elo1978rating}. To help the Preference RM penalise off-topic answers, we add a randomly chosen \emph{distractor} response to each comparison, sampled from the rest of our response preference data. We also found that two auxiliary losses improved preference modelling. We add a classification loss predicting whether evidence conditioned answers were supported and plausible, following \citep{menick2022teaching}. We also ask raters to indicate when all responses in a comparison are low quality and regularise the corresponding Elo scores to be negative. Refer to \ref{app:reward-models} to see how auxiliary losses from these tasks are incorporated, and how Chinchilla was fine-tuned for this task.

The Rule RM is a conditional classifier $r(x, y) \in [0, 1]$ that estimates the probability that the rule $y$ was violated by Sparrow at any point in the dialogue $x$. Rule RMs are trained on rule violation data (\ref{itm:adversarial-probing}). 
We use a version of instruction tuning~\citep{gao2020making,wei2021finetuned,saeidi2021cross,kotonya2022policy} as we find it gives good performance with small amounts of data (see \ref{sec:general}). The training objective is to maximise the likelihood of the sequence of tokens corresponding to \emph{Yes} or \emph{No}, depending on the label from human ratings, given the prompt in \ref{fig:rule_classifier_prompt} formatted with the corresponding dialogue and rule. Because the Rule RM is trained jointly on all rules, memory and computation can be shared across rules for the same dialogue, such that memory and computation scale weakly with the number of rules; refer to \ref{app:reward-models} for details.

In all cases when fine-tuning, we freeze the bottom 64 transformer layers of Chinchilla, and only fine-tune the final 16 layers; this allows sharing of the frozen layers between the rule model, preference models, and the base LM/policy when reranking and during reinforcement learning training, resulting in a reduced memory footprint (\ref{fig:hydra-diagram}).

\subsection{Reranking} \label{sec:rerank}

Given a Preference RM and a Rule RM, a dialogue agent's policy can be improved by reranking multiple sampled responses as in \citet{thoppilan2022lamda, askell2021general, menick2022teaching}.  At inference time, we draw $N$ samples and select the sample with the maximum combined reward. We call such models \textit{`model@N'}. 
\ref{fig:reranking} shows inference time operation of Sparrow with reranking \emph{@8}. Given the previous dialogue, a generative model samples four answers using a standard dialogue prompt (\ref{app:prompt}) and two search queries using an evidence prompt (\ref{app:prompt_evidence}). The search queries are used to retrieve up to four search result fragments, which in turn are used to sample Sparrow responses (with the fragments shown expressed as evidence). The total of 8 samples are rescored according to \ref{eq:reranking}, in a scheme loosely inspired by the product of experts approach \citep{productofexperts}. Here $R_{pr}$ is the Preference RM score, $AVG(R_{pr})$ is the average Preference RM score on the valid set, and $R_{rule_i}$ is the Reward RM score of rule $i$ out of $n$ (the probability of the rule being followed, so that higher is better).

\begin{equation}
\label{eq:reranking}
  R_{\text{rerank}} = \frac{e^{R_{pr}}}{e^{R_{pr}} + e^{AVG(R_{pr})}}  \left( \prod _{i=1}^{n}R_{\text{rule}_{i}} \right)^{\frac{1}{n}}
\end{equation}

Reranking also enables our agent to decide whether to make use of search results and provide evidence. This ability can be viewed as a selective prediction of using evidence (or prediction with {\it a reject option}) \citep{el2010foundations, geifman2017selective, geifman2019selectivenet, kamath2020selective}. The preference RM gives high scores to factual model responses with clearly supporting evidence and responses without evidence to non-factual questions. It gives lower scores for responses with unnecessary or low-quality evidence. The Rule RM penalises responses that break rules.

\subsection{Supervised fine-tuning} \label{sec:sft}

Supervised fine-tuning (SFT) via LM loss is the main training technique used by LaMDA~\citep{thoppilan2022lamda} while the Anthropic assistant~\citep{bai2022training} instead uses \emph{context distillation}, and otherwise relies on reward modelling and reinforcement learning. We also fine-tune Chinchilla directly via LM loss on the collected dialogues rated as preferred and rule compliant, as an alternative to reward modelling and reinforcement learning. For per-turn preference data, we fine-tune the model to produce the preferred response. For adversarial probing dialogues, we fine-tune the model on the \texttt{Agent} responses in dialogues rated at least \emph{good} (\ref{sec:eval_start}) and where no rule was broken. The SFT model provides a stronger baseline than DPC, as well as a better initial starting point for RL.

\subsection{Reinforcement learning} \label{sec:rl}

Similar to \citep{bai2022training}, we use reinforcement learning (RL) with our reward models to improve the dialogue agent. This approach complements reranking, which is expensive at inference time; RL is expensive to train but adds no inference cost, and the two can be combined freely.

Our RL scheme is illustrated in \ref{fig:rl-system}. Each episode consists of a single statement (not a complete conversation) conditioned on a preceding dialogue context, where the actions are individual tokens and the reward is given at the end of each episode (\ref{app:rl-rewards}).

Unlike \citet{bai2022training} who perform RL on single-statement continuations of previously collected human-agent dialogues, we use a form of self-play, where during training the generated statement and the dialogue context form a new dialogue context for a later episode; thus, Sparrow generates multiple turns of a dialogue, playing the role of \texttt{User}, \texttt{Agent}, and \texttt{Search Query} (\texttt{Search Results} are retrieved programmatically) over multiple episodes. Note that \texttt{Search Query} statements are treated as separate episodes from \texttt{Agent} statements. For each episode, the preceding dialogue context is prefixed with a prompt specific to the role Sparrow is playing in that episode (\ref{app:rl-prompts}). The preceding dialogue context can come from several possible sources, which are effectively \emph{user models} that exhibit different interests and behaviours:

\begin{itemize}
    \item \textbf{A dataset of questions.} We use the filtered train subset of ELI{5} from GopherCite \citep{eli5, menick2022teaching}.
    \item \textbf{A conversation with a human.} We take a mixture of open-ended and adversarial conversations from annotators and randomly truncate them to allow Sparrow to continue the conversation from an intermediate turn.
    \item \label{itm:red-teaming} \textbf{A red team language model.} We use the zero-shot method of \citet{perez2022redteaming} by prompting Chinchilla to generate adversarial questions that augment the available human data (\ref{app:red-teaming} details these prompts). 
    \item \textbf{Self-play data accumulated through training}. During training, Sparrow generates a response to each dialogue context in a batch, playing the role of both \texttt{User} and \texttt{Agent} as needed. Any valid statements (\ref{app:rl-rewards}) are combined with their dialogue contexts to form a new context that is added to a self-play buffer, up to a maximum conversation length of 12 statements. This allows Sparrow to learn by talking to itself.
\end{itemize}

This amounts to optimising the RL policy conditioned on a distribution of conversational contexts induced by the above mixture. That is, the optimisation objective is

$$
\arg\max_\pi \mathbb{E}_{c\sim \mathcal{D},\ s\sim \pi}\left[R(s\lvert c)\right],
$$

where $c\sim\mathcal{D}$ is a distribution of dialogue contexts defined above, and the $s = a_{1:T}$ are utterances generated according to the agent's policy $\pi$. Note that we elide the summation of rewards over the episode as the reward is zero at all steps apart from the end of an episode, and we don't apply explicit discounting. The reward function $R$ is defined in full in \ref{app:rl-rewards}.

All statements after the initial dialogue context are generated by Sparrow, taking the role of \texttt{User}, \texttt{Agent}, or \texttt{Search Query} as needed. Future work could extend this to a league of user models optimised to probe different aspects of the main agent's behaviour \citep{alphastar}.

The RL reward is given by the sum of the response preference and rule violation models, where the rule reward is the mean over all rules scores, combined with programmatic rewards for validity and conciseness (see \ref{app:rl-rewards}). \texttt{User} statements do not receive rule rewards, but are trained by the same preference model as \texttt{Agent} statements. Due to the different output ranges of the preference and rule models, we independently normalise each one using a running mean and standard deviation before adding them.

The dialogue context, sampled actions, and rewards from the trajectory data are used to update the model parameters. The RL algorithm we use is a batched synchronous advantage actor-critic~(A2C; \cite{mnih2016a3c}), or equivalently REINFORCE with baseline~\citep{Sutton1998}; we found that V-MPO \citep{song2019v} did not improve performance significantly and is computationally more expensive. Due to nucleus sampling, our training data is off-policy, which we do not correct for; one solution could be to introduce off-policy methods.

We initialise the policy to either Chinchilla or an SFT model (\ref{sec:sft}); Sparrow was initialised to the SFT model at RL training time. To prevent RL from collapsing to a single, high-reward generation, we penalise the KL divergence between the fine-tuned policy and the initial \emph{teacher} language model. To mitigate the memory requirements for multiple Chinchilla-sized models --- multiple reward models, policy, value, and teacher models,  which must all fit in device memory --- we train only the top layers of each and fuse them into a multi-headed \emph{hydra} model, with a separately trained `head' for each model and a shared trunk of pretrained parameters (\ref{fig:hydra-diagram}).

The use of self-play, search, fine-grained rules, and LM red-teaming extend beyond the proposals of \citet{bai2022training}.  \ref{fig:rl-ablations} explores the impact of rules and red-teaming in more detail, showing that introducing red-teaming data during training is complementary to the use of rule models. Varying the data distribution together with rewards is an expressive means for shaping behaviour, and we consider it under-explored in the current version of Sparrow. A long-term approach should make the trade-off of helpfulness and harmlessness test-time configurable \citep{abdolmaleki2020distributionalmompo} and train over an expanding universe of trade-offs and topics in an open ended fashion \citep{xland} to find an optimal training data distribution.

\section{Results and analysis} \label{sec:results}

\subsection{Preferences and rule violations} \label{sec:headlines}

Our primary evaluations for information-seeking dialogue, shown in \ref{fig:pareto_plot}, are conducted by asking paid annotators to assess model responses in two types of human data collection: per-turn response preference and adversarial probing (\ref{sec:eval_start}). In both cases, the evaluated models are shown to the individual raters in a round-robin fashion.

\paragraph{Three-model preference rate} We assess the quality of a model’s answers in terms of preference against two DPC baselines. \textit{DPC - never search} is a prompted model without search (\ref{app:prompt}). \textit{DPC - always search} is a prompted model that is forced to produce both search query and search results at every model turn (\ref{app:prompt_evidence}). All evaluated models are able to select whether to search and provide evidence. We use three-model comparisons rather than pairwise preference to avoid biases causing the raters to default to preferring the option with or without evidence without careful evaluation. The three-model preference rate is established through per-turn preference comparison of an evaluated model with the two DPC baselines. Each dialogue task starts with a \texttt{User} turn sampled from a test set of 200 utterances, consisting of 100 randomly sampled questions from the ELI5 dataset~\citep{eli5} (filtered for toxic content), and 100 sampled from free dialogues with annotators who were instructed to ask Sparrow factual questions.

\paragraph{Violation rate under adversarial probing} We ask the raters to lead a conversation with Sparrow in such a way that Sparrow might break the specified rule (one of first 18 rules in \ref{tab:harm-rules}) as described in \ref{sec:eval_start}. We aggregate by dropping \emph{unsure} ratings and binarising the scale into  \emph{break} and \emph{follow}.

Optimising for preference and harmlessness stand in opposition to each other~\citep{askell2021general}. For example, an agent that always responds with ``I can't answer that'' is perfectly harmless but not very useful, while an agent that always engages with the question may be led astray by malicious users and emit harmful language. To express this trade-off, we present our evaluations in the form of a Pareto frontier in \ref{fig:pareto_plot}. Of all models, we find that combining RL with \emph{reranking@8} (in red) achieves the best performance both in terms of preference win rates and resilience to adversarial probing.

RL and reranking are complementary: \ref{fig:reranking_improves_win} shows that reranking gives a consistent three-model preference rate improvement for all the classes of models (DPC, SFT, RL). \ref{fig:model_improves_violations} shows that RL and SFT outperform the DPC baseline by having lower violation rates under adversarial probing.

\ref{fig:violations_per_rule} shows that our interventions improve Sparrow's resilience to attack for a majority of rules. However, they do not alleviate harms from the following rules: \emph{no stereotypes}, \emph{no medical advice}, \emph{no legal advice}, \emph{no microaggressions}, and \emph{no insults} (please refer to \ref{app:error_cases} for examples of successful and avoided attacks). We hypothesise that this is caused by the following factors:

\begin{itemize}
    \item Sparrow often finds convincing search results supporting the responses for medical or financial topics, or even stereotyping opinions from the web (we do not block forums).
    \item Due to rater well-being concerns, we collected less data for some rules. All the above-mentioned rules (\ref{app:harm-rules}) fall into that category. \ref{tab:RuleRMDataset} shows data collected per rule.
    \item Many of the human raters for the Preference RM data have never completed the \emph{adversarial probing} or \emph{rule rating} task and so may unknowingly pick rule-breaking responses.
\end{itemize}

\subsection{Evidence evaluation} \label{sec:evidence-eval}

\paragraph{Multi-turn supported and plausible evaluation} We assess Sparrow's responses and accompanying evidence through human evaluation, using the metrics of \emph{supported} and \emph{plausible} as defined in \ref{itm:evidence-app} and GopherCite \citep{menick2022teaching}. We evaluate these metrics in the multi-turn dialogue setting as an extra rating task (\ref{itm:evidence-app}) within the per-turn preferred response task (\ref{sec:eval_start}). We measure the supported and plausible rates achieved on the turns requiring factual responses from the model (as determined by raters). \ref{tab:s_and_p} shows the rate at which individual models chose to provide answers with evidence, along with the supported and plausible rater judgements for the cases in which the evidence was given. We find that humans determine our best model's responses with evidence to be plausible and supported in 78\% of the cases.

\paragraph{Selective prediction of using evidence} An important ability of the agent is to determine for which turns to display supporting evidence alongside the response. Sparrow should not condition on and show evidence for responses to questions such as \textit{``How are you?''} or when evidence would lead to rule violations; however, it should search and provide evidence for factual questions like \textit{``What is the radius of Earth?''}. We evaluate this ability with the annotation tasks described in \ref{itm:evidence-app}: given the previous dialogue ending with a \texttt{User} turn, the rater indicates if the \texttt{Agent} turn requires grounding in external knowledge. Since our primary test set consists mostly of information-seeking dialogue conversations, we additionally include 100 conversational questions; these were generated by Chinchilla by asking it for a list of \emph{questions to ask someone} (\ref{app:red-teaming}).
The confusion matrix in \ref{fig:confusion_matrix} shows that Sparrow generally agrees with raters on whether evidence is necessary, with an overall agreement rate of  over 90\%.
We find this to be a particularly strong result, given that we only used per-turn preference data for training.

\paragraph{False negatives} We were particularly interested in the 7\% of cases where raters judged that external evidence should be cited, but Sparrow did not (marked with an asterisk in \ref{fig:confusion_matrix}). 51\% of the time, raters actually changed their minds after seeing Sparrow's response and agreed that evidence would not be useful. Qualitatively, we found three common explanations for the remaining half: a) questions whose answers would normally require evidence but which would lead to rule violations (e.g. medical questions) and where Sparrow (correctly) declined to answer, b) cases where all the search \emph{results} were of low quality, and so reranking picked a non-search response, and finally c) simple mislabelling by the raters.

\paragraph{Comparison to GopherCite} Sparrow's ability to support its responses with evidence extends the methods of GopherCite~\citep{menick2022teaching} to the interactive dialogue setting. GopherCite was designed for single-turn question answering and does not generalise to dialogues with followup questions. Given these differences, we compare GopherCite to an \emph{always search} Sparrow which only considers answers with evidence during reranking. We evaluate Sparrow with reranking over 4 responses with evidence (\emph{RL@4 - always search}), and GopherCite with reranking over 16 responses as in~\citep{menick2022teaching}.

We compare GopherCite to Sparrow head-to-head in the question answering setting, using the GopherCite human evaluation interface and test set (FilteredELi5). In \ref{tab:s_and_p_eli5_gophercite} we find that in this setting Sparrow (\emph{RL@4 - always search}) achieves similar supported and plausible rates to GopherCite. Human raters also show a preference 63\% (90\% CI=[56\%, 70\%]) for Sparrow answers over \emph{GopherCite RL@16} when comparing model responses in this setting. These results show that Sparrow, an interactive system that can additionally answer follow-up questions in real-time, does not degrade QA performance as compared to the larger and slower GopherCite system.

\subsection{Correctness evaluation} \label{sec:author-eval}

It is naturally of interest how often Sparrow is correct during a conversation. However, robustly assessing correctness in an open-ended setting is challenging. Our supported and plausible evaluations do not require human raters to make an absolute judgement of the response correctness or to fact-check with external sources, instead only asking if a response is supported and plausible given the model-provided evidence. Such statements are not necessarily factually correct (\ref{sec:limitations}). In addition, supportedness evaluations are not possible for model statements without evidence.

To give a coarse notion of correctness, we carried out an additional small-scale investigation. We collected 200 information-seeking dialogues instructing raters to ask factual questions and follow-ups. In this ``free dialogue'' setting, participants were not instructed to probe for rule violations, or briefed on the rules the model should follow. Of these dialogues, 100 conversations were collected from the baseline DPC without evidence, and 100 were collected from Sparrow (\emph{RL@8}).

These dialogues were then annotated by some of the authors for correctness, according to the following procedure:

\begin{enumerate}
\item Rate just the model response, ignoring any evidence. Rate the correctness of each claim based on general knowledge and fact-checking with external sources. Assign scores on a Likert scale of: \textit{false, mostly false, unsure, mostly true, true}. If the last turn requires no externally-verifiable claims (small talk or questions about Sparrow itself), rate the turn as \textit{not applicable}.

\item Rate the evidence if present. Determine whether the evidence is helpful and sufficient to verify the correctness of the model response. Assign a rating according to a Likert scale of: \textit{not supportive/irrelevant, mostly unsupportive/irrelevant, unsure, mostly supportive, supportive}
\end{enumerate}

We release the transcripts and our ratings: \url{https://dpmd.ai/sparrow-samples}

We do not judge the model responses for helpfulness (e.g.\ properly answering the question), only for correctness of factual claims. To aggregate correctness judgements, we drop each \textit{not applicable} or \textit{unsure} and binarise the Likert scale.

\ref{tab:correctness} shows the results of this investigation, providing some evidence that conditioning answers on the search results and training for quality of evidence increases the correctness of Sparrow's answers. We stress that this evaluation for correctness is not on an adversarial distribution. If you aim to find cases where Sparrow is incorrect, it is not challenging to do so, and some examples are included in \ref{app:interesting-samples}. However for the questions our raters chose to ask in the free conversation setting, Sparrow gave answers that were rated as true or mostly true around 80\% of the time in our judgement.

\subsection{Rater perception of trustworthiness} \label{sec:rater-perception-eval}

After each free dialogue was collected for the correctness evaluation in \ref{sec:author-eval}, we asked the raters \textit{How trustworthy did you find the AI?}, on a scale of \emph{very trustworthy, trustworthy, somewhat trustworthy, neutral, somewhat untrustworthy, untrustworthy, very untrustworthy}.

This is a subjective measure reflecting the personal biases, varying expectations, and conversational tendencies of individual participants. We reduce the impact of some confounders by presenting the models to participants in round robin fashion, so that each participant interacts with both models (starting from a random one), and also by setting a minimum response time of 6 s in order to avoid effects driven by latency.
\ref{tab:overall_and_trustworthiness} shows that participants give higher ratings to Sparrow than the DPC baseline for trustworthiness, indicating that they find Sparrow trustworthy or very trustworthy in around 70\% of conversations.

\subsection{General versus\ specific rules} \label{sec:general}

The rules defined in \ref{sec:rules} are a key part of the adversarial probing and rule rating evaluations described in \ref{itm:adversarial-probing}. While they vary in granularity, our rules are all fairly specific. Prior work has either used a broad definition of harm or trained a single reward model for all harms~\citep{bai2022training, thoppilan2022lamda}. Here, we analyse the effectiveness of our fine-grained approach compared to using only a single \emph{general} harm rule: \emph{Do not make statements that could be considered harmful (by harming the user or otherwise).} We focus on harm rules since harmful language can be a long-tail phenomenon for which adversarial probing is particularly important \citep{zhang2019hate}.

To assess how effective specific rules are compared to the general harm rule, we collected 1078 dialogues with DPC. For 538, participants were instructed to adversarially probe for violations of a specific rule; the other 540 were asked to adversarially probe for violations of the general harm rule. In the latter case, the specific harm rules were listed in the instructions as in \citet{thoppilan2022lamda}. For all of these dialogues, we randomly sample from two \texttt{Agent} prompts, the DPC prompt in \ref{tab:DPCprompt} and the less safe prompt in \ref{tab:lesssafeprompt} which makes the \texttt{Agent} more vulnerable to violations. All of these dialogues were then independently re-rated against all rules, including the general harm rule. Each rater evaluated at most 5 rules per dialogue to avoid fatigue in the re-annotation phase and each dialogue was rated for each rule by 2 raters independently. Re-annotating all conversations for all rules is necessary for this comparison, but is not our usual protocol.

\paragraph{Effectiveness of adversarial probing} \phantomsection\label{itm:general_adversarial}
To train a rule model with high accuracy for many rules, the training data needs to sufficiently cover the space of harms.  \ref{fig:general_violations} shows that adversarial probing for a specific rule lets us steer raters towards problems that we lack data on. If raters are asked to target a specific rule, they are more likely to elicit a violation of that rule than if the raters are probing for the general harm rule. This effect is particularly notable for rules like \emph{do not offer financial advice}, which raters seem less likely to think of when probing (despite all rules being listed in the instructions as examples of harm).

\paragraph{The general harm rule as a method to find new specific rules} \phantomsection\label{itm:general_novel}
By definition, specific harm rules cannot cover the entire space of harm. A general harm rule might act as a catch-all to find and fix bad behaviour not covered by specific rules. Indeed, we find that at least 19 of 566 dialogues that adversarially probed the general harm rule discover novel harms not covered by our specific harm rules. The discovered novel harms all fell under the Information Hazards and Misinformation Harms categories described in \cite{weidinger2021ethical}. See \ref{app:uncovered-harms} for more details.

\paragraph{Effectiveness of rule rating} \phantomsection\label{itm:general_rule_rating}
We investigate how using specific rules impacts inter-annotator agreement (IAA) compared to using a general rule. The IAA is computed as Krippendorff's Alpha~\citep{krippendorff2011computing}, by binarising the Likert scale for rule violations into \emph{break} and \emph{follow}, discarding \emph{unsure} ratings. To compare on the same task, we merge the specific rule annotations for any given text into a single \emph{was any specific harm rule violated} rating. The IAA for the \emph{was any specific harm rule violated} rating is 0.53 (95\% CI=[0.47, 0.59]), while it is 0.37 (95\% CI=[0.29, 0.46]) for the \emph{was the general harm rule violated} rating for the same dialogues; indicating higher IAA when asking about specific harms rather than general harmfulness. See \ref{fig:general_iaa} for per-rule IAAs. 

\paragraph{General versus\ rule-conditioned rule reward model} \phantomsection\label{itm:general_rule_classifier}

Rule-conditioned RMs perform better compared to general safety classifiers (as used by \citet{thoppilan2022lamda, xu-etal-2021-detoxifying}), while requiring less data. For this comparison, we use the same set of $N=538$ rule-specific adversarial dialogues and their rule-specific annotations (excluding general harm annotations), split into 30\% test and 70\% training data. If at least one of the two raters considers the rule to be violated, we also consider it violated. When training the overall safety classifier, we mark each dialogue as unsafe if any rule is violated. This allows us to train both classifiers on the same data (rule-specific annotations), test on the same task \emph{was any rule violated}, and compare head to head. For the rule-conditioned classifiers at test-time, we predict \emph{was any rule violated} as the maximum rule-conditional violation probability across all rules. We trained both a Rule RM using instruction tuning, and an unconditional general rule classifier using a linear head on the full training set, and independently on 50\% of the full training data (randomly sampled) to investigate sample efficiency; both models were optimised by sweeping over the same hyper-parameters (see \ref{app:reward-models}). \ref{fig:rule_classifier} shows that rule conditioned classifiers achieve a higher final performance (0.85 vs 0.77 AUC) when trained on the full training set, as well as better sample efficiency (0.81 vs 0.69 AUC) on 50\% of the training data (evaluated on the same \emph{was any rule violated} test set). Refer to \ref{app:reward-models} for details on Rule RM training.

\subsection{Distributional harms} \label{sec:distributional}

Targeted rules and inline evidence mitigate instance harms, which can be identified given a single example, but not distributional harms, which depend on aggregate behavior. For example, if Sparrow correctly follows our rule about stereotypes\footnote{Our stereotype rule is ``Do not use stereotypes or make any other harmful generalising statements about groups of people.''  Please see \autoref{tab:harm-rules} for a list of all our rules.}, it will not make statements like ``women don’t make good scientists''. However, it could still exhibit a bias in aggregate by rarely discussing women when asked about notable scientists. Prior work outside dialogue has shown that mitigating instance harms, like toxicity, can exacerbate distributional harms~\citep{welbl-etal-2021-challenges-detoxifying,xu-etal-2021-detoxifying}.

\citet{blenderbot3} and \citet{bai2022training} also analyse distributional harms in dialogue models, and we extend their work in several ways. In \ref{sec:stereotypes}, we show that bias in dialogue models can be amplified after multi-objective RL, that bias persists in model samples (likelihood-only metrics may compare two unlikely continuations), and that bias can manifest in incorrect answers on fully specified questions (in addition to bias in word correlations as shown in \cite{blenderbot3}). In \ref{sec:disparate-impact} we conduct a preliminary investigation into disparate impact (differences in model performance for different groups), focusing on fact-based question answering. 

\subsubsection{Stereotypes and social biases} \label{sec:stereotypes}

We first consider datasets which test whether models rely on harmful stereotypes. Models which favor stereotype-reinforcing answers can cause harm in two ways: they may reinforce stereotypes, and they may falsely characterise individuals in stereotyped groups~\citep{parrish2021bbq}.

\paragraph{Setup}
We use three datasets designed to test models' reliance on stereotypes: Winogender \citep{rudinger2018gender}, Winobias \citep{zhao2018gender}, and BBQ \citep{parrish2021bbq}. Winogender and Winobias are co-reference datasets, with Winobias including two sentence ``types''; \emph{type 1} sentences are designed to be more challenging as they do not include syntactic cues for coreference. BBQ is a question answering dataset which asks questions about people belonging to different groups based on context provided in a few sentences. In each dataset, questions can be answered with either a stereotype-reinforcing or stereotype-challenging response (and on BBQ, an additional ``I don't know'' response). For Winogender and Winobias, we follow \citet{rae2021gopher,gpt3,hoffmann2022training} and select an option by comparing LM likelihoods given a zero-shot dialogue prompt. For BBQ, we instead follow \citet{parrish2021bbq} and sample responses. This directly measures bias in the LM outputs, and avoids comparing low likelihood continuations to each other. We use a 5-shot dialogue prompt to ensure the model uses the expected output format.

For our bias metric $s$, we measure the fraction of stereotype-reinforcing versus stereotype-challenging responses, as proposed by \citet{parrish2021bbq}. We rescale so that $s=1$ indicates always being stereotype-reinforcing, $s=-1$ always stereotype-challenging, and $s=0$ an even balance. $s$ is also the difference in accuracy between questions with stereotype-reinforcing versus stereotype-challenging answers, and a perfectly accurate model would have $s=0$ (see \ref{app:stereotypes}).
For BBQ when ``I don't know'' is correct, we follow ~\citet{parrish2021bbq} and rescale the bias score (defined as $s_{\textnormal{ambig}}$ in \autoref{app:stereotypes}) to reflect that a model which correctly abstains from answering questions is preferable.
\ref{app:stereotypes} has full details on our datasets, metrics and setup.

\paragraph{Results}
\ref{fig:winogender_bbq} shows our results. We find that bias persists across models and datasets.  On Winobias type 1 questions, both the DPC and RL models are roughly 36\% (absolute) more likely to be correct when it is stereotype-reinforcing. RL finetuning can amplify bias over the base model: on Winogender, the bias score increases from 0.06 to 0.10. For ambiguous questions in BBQ, bias scores increase in 10 out of 11 categories. Averaged across groups, the bias score increases from an average of .05 to 0.10, with larger effects in some categories such as physical appearance, disability status, and age. Evidence in \ref{app:stereotypes} suggests much of this effect is due to the RL model becoming less likely to abstain, along with a tendency towards stereotype-reinforcing responses in such cases.

\subsubsection{Disparate impact for factual question answering} \label{sec:disparate-impact}

Disparate impact might arise if our system is less useful for different groups. Here, we aim to more directly study how disparate impact might arise in an information-seeking task by measuring our model's ability to answer questions about specific groups. Though this does not directly measure usefulness for different groups (which is more difficult to do), it may be correlated, and also provides practice in aiming towards systems which benefits all users equally.

\paragraph{Setup}
Following \citet{gor2021deconf}, we evaluate factual question answer performance across questions relating to different demographic groups (gender, country, and occupation) using three QA datasets (Natural Questions \citep{kwiatkowski2019natural}, Quiz Bowl \citep{graber2012besting}  and TriviaQA \citep{joshi2017triviaqa}). We give questions directly to the dialogue model and report the rate at which the correct answer appears within the model's response (\emph{exact match} accuracy) for each group.

\definecolor{evcolor}{RGB}{173, 112, 255}
\definecolor{rlcolor}{RGB}{18, 54, 147}

\paragraph{Results}
Given the task's emphasis on facts, we observe the largest effect sizes from incorporation of evidence. We thus focus on these effects, leaving full results to \ref{app:disparate}.
\ref{fig:deconf-tqa} shows results for the largest dataset, TriviaQA, where incorporating evidence improves accuracy across all categories. 
\ref{tab:p-values} reports when correlation between accuracy and demographic group is significant, per a $\chi^2$ test.  Similar to \citet{gor2021deconf}, we do not always see a statistically significant effect, and including evidence can both introduce and remove correlations.

\section{Discussion}
\label{sec:discussion}

As discussed in \ref{sec:intro}, we view this paper as a base on which to build and investigate further safety mitigations.  There are several major directions we hope to explore going forwards.

\subsection{Evidence limitations} \label{sec:limitations}

A key limitation of Sparrow is that we use only one external knowledge fragment at a time, in contrast to WebGPT~\citep{nakano2021webgpt} and LaMDA~\citep{thoppilan2022lamda}. WebGPT also allows scrolling within retrieved pages and clicking on links.  SeeKeR~\citep{shuster2022seeker,adolphs2021reason} uses an intermediate knowledge-extraction step to reduce from several documents concatenated together to a smaller \emph{knowledge sentence}, while we rely on search engine text snippets. Our previous work \citet{menick2022teaching} selected an evidence fragment from a longer context, a feature which we removed due to pressure on the model's context length as Sparrow's context holds an entire dialogue history and a prompt.  We believe these limitations are best addressed via multistep reasoning~\citep{lewkowycz2022solving,creswell2022selection,dohan2022language}, with the further benefit of producing interpretable reasoning traces. Sparrow also frequently copies text verbatim from the evidence (\ref{fig:sample_searborn,fig:css-sample} are examples), which could likely be mitigated with further rules.

In this work, we roughly say a statement is correct if each factual claim it makes is supported by either common sense or evidence from a source that is trustworthy. This breakdown rules out some true statements, but is conservative and supports evaluation by human raters; see \citet{evans2021truthful} discussion. Showing this evidence also to downstream users gives agency in choosing whether to trust model statements. However, we do not investigate the trustworthiness of sources in this paper, and this breakdown does not account for statistical evidence such as aggregating many different sources together. Finally, although we believe RLHF and evidence are key for correctness, other machinery such as interpretability~\citep{elhage2021mathematical} or eliciting latent knowledge~\citep{christiano2021elk} will be required to specifically target \emph{honesty} as advocated by~\citet{askell2021general}.

\subsection{Dialogue as a supervision mechanism}
\label{sec:dialogue-supervision}
In this paper dialogue is the task, but our long-term hypothesis is that dialogue is a key component of accurate supervision for machine learning models. Indeed, we chose dialogue as the task in part to build experience and infrastructure to tackle dialogue for supervision. Prior work has suggested that assistance from ML systems may help with accurate human supervision~\citep{irving2018debate,christiano2018amplification,leike2018scalable}. In such cases, dialogue is a natural medium for this ML assistance, as it enables both clarifications and iterative discussion of subtle points. Determining whether a particular model behaviour is good is often quite subtle, and it is easy for human review (whether by paid crowdworkers or the authors of this paper) to miss key details or misinterpret text.

\ref{fig:subtle-dialogue} shows an example of a dialogue about whether a model is violating a rule, in this case edited from a Slack discussion by several of the authors. In this case, as our task is also dialogue, the supervision dialogue is about a dialogue transcript, but one could also have a supervision dialogue about non-dialogue behaviour (e.g., a generated image).  The initial statement (by an author of the paper) is incorrect, someone else provides a correction, and the first person changes their mind.  But then another author points out a different flaw. The eventual conclusion is that the first rule is not violated, but a different rule might be.

Our hypothesis is that this type of multistep discussion is required to resolve subtle cases of supervision correctly. In the above dialogue, humans provided the corrections and clarifications, but sufficiently capable dialogue agents could also provide them. The same principle applies with cited evidence, as additional sources or arguments may be needed if an initial source quotation is taken out of context. The adversarial case of dialogue for supervision is debate, where two or more dialogue agents point out flaws in each other's statements~\citep{irving2018debate}. However, dialogue for supervision also needs cooperation between humans and agents to jointly clarify what is meant, and avoid misunderstandings or gaps~\citep{hadfield2016cooperative,russell2020compatible}. Determining the best way to combine adversarial and cooperative behaviour will be key as we move towards dialogue for supervision. Initial work towards multistep human interaction methods includes simulated debate using frozen question answering models~\citep{perez2019finding} and recursively summarising books \citep{wu2021recursively}, which simplifies the rating task from evaluating book-length summaries to passage-length summaries. Initial evidence from one-step debate is mixed: \citet{saunders2022self} find that model-generated critiques help humans notice flaws in summaries, but in \citet{parrish2022single} accuracy did not improve when humans were shown explanations.

\subsection{Ethical and sociotechnical aspects}
\label{sec:external-inputs}

A primary goal of the rule mechanism is to enable the scalable incorporation of input from multiple stakeholders --- including users and affected groups --- on what constitutes good speech for language agents. However, the successful implementation of such a mechanism raises a range of open research questions. For example, any rule mechanism will need to consider the origins of its rules and balance the needs and expectations of relevant stakeholders. In this study, the rules were generated in consultation with domain and legal experts and centered around a small set of proposed rules. In future, more participatory inputs~\citep{lee2019webuildai,halfaker2020ores,berditchevskaia2021participatory} from other stakeholders will be critical for developing language agents that are both legitimate and aligned to the needs of its users. Participatory approaches are not a panacea, however, and their successful deployment turns on a set of technical and ethical considerations that have been well documented in prior research on sociotechnical ML~\citep{birhane2022power,sloane2020participation}.

We distinguish two goals of rules in influencing agent behaviour: mitigating harms and incentivising better speech. Prior research from \citet{bender2021dangers} and \citet{weidinger2021ethical} has delineated a range of emergent and existing harms from large language models, and \citet{rauh2022characteristics} describes six characteristics along which language harms can vary, including some specific to dialogue. The impact of these harms is not distributed evenly, as underrepresented groups are most likely to be at risk due to problematic agent behaviour~\citet{tomasev2021fairness}. We can also use rules to incentivise speech that is more closely aligned with appropriate norms and values: \citet{kasirzadeh2022ideal} build on work by \citet{grice1975logic} in formulating \emph{pragmatics} principles whose joint enforcement results in effective and beneficial communication.  Using rules to shape dialogue can be important both for dialogue as a task and dialogue for supervision, where our goal is the accurate evaluation of agent behaviour. Pragmatics may be crucial when using dialogue to supervise highly capable agents: there are many types of deceptive argument to detect~\citep{schopenhauer1831art}, and these may differ from normal human-to-human communication~\citep{irving2019social}.

The existence of a potentially large number of rules motivates techniques which scale to many rules. Our rule-conditional reward models work well up to the number of rules used in this paper, but we expect further architectural work to be required to scale to 100s or 1000s of rules. Finally, a key practical advantage of collecting data via detailed rules is that conflicts and weighting between rules can be changed after the fact: \citet{saeidi2021cross} express policies as expression trees with rules as the leaves, with the expression either written by experts or inferred from prose~\citep{kotonya2022policy}.

\subsection{More cognitive science research is needed} \label{sub:cognitive-science}

Since our goal is to help humans supervise dialogue agents, understanding whether we have succeeded at our task depends fundamentally upon insights from cognitive science and human computer interaction~\citep{irving2019social}. This analysis is particularly important for interactive settings such as dialogue with complex interdependencies between agent responses and human beliefs and preferences. Here we discuss two important topics for future research; there are many others.

First, a core goal in our research and others is to ground agent responses in evidence~\citep{evans2021truthful}. While this is a critical antidote to harms arising from false or misleading statements, treating truth and evidence only as a property of model outputs misses downstream effects on the minds of the human conversational partners. Extensive literature demonstrates that strong beliefs can resist change despite compelling contradictory evidence~\citep{gershman2019}. Numerous mechanisms for this have been proposed, the most well-known of which is that of the motivated reasoning bias~\citep{kunda1990}. Finding modes of evidence that are less susceptible to such cognitive biases will be important for the future of aligned AI and beneficial human-AI interaction.

Second, as the space of potential rules to apply increases, we must ask which granularity is most appropriate. It is usually possible to find increasingly granular, specific rules in any given category of harm. Intuitively, more specific rules seem easier for human raters to apply, but a single human will be unable to hold in mind more than a handful of rules at a time (we limit our own evaluations to at most 5 simultaneously). There is therefore a trade-off between rule specificity and efficiency in the data collection. In principle, this is a question that can be addressed empirically with suitable human experiments.

\subsection{Broader impacts} \label{sec:impacts}

As discussed in section 7.3 of \citet{rae2021gopher}, we believe most language harms are best mitigated downstream of LLM pretraining, due to faster iteration cycles, application-dependence of harms, and multiple roles served by a single model (we use Chinchilla as both policy and classifier).  This work is one component of this downstream mitigation, but our methods are limited to instance harms detectable by raters without significant help.  Issues such as privacy~\citep{abadi2016privacy} and social, linguistic or environmental justice~\citep{bender2021dangers} require mitigations at pretraining time in addition to downstream work, though rules have a role (such as teaching an agent to not reveal information that should be private, even if it is available on the open web).

Like many alignment methods, ours are dual-use: they could be used to enforce harmful rules as easily as beneficial ones. To avoid harmful outcomes we must address how control over the rules is decided, whether affected parties share in this control, and whether they have visibility into what rules are in effect; considerations analogous to those raised by \citet{denton2020people} for datasets.

\section{Conclusion}

Building helpful, correct, and harmless agents out of raw generative models involves both \emph{width} and \emph{depth}: width to deal with the detailed complexity of goals and topics, and depth to handle each of these carefully and correctly.  With Sparrow, we have focused on width: breaking down goals into detailed rules, and allowing the agent to pull in external knowledge to broaden the topics it can correctly discuss. We found that these techniques work, enabling Sparrow to respond helpfully more often as measured by rater preference, correctly cite evidence 78\% of the time for factual questions, and reduce rule violation rate to 8\% under adversarial conditions. Addressing depth will require multistep reasoning for the agent to talk through problems with itself (leaving interpretable traces for humans to evaluate), expert and participatory engagement to find and evolve good sets of rules, debate and dialogue for supervision to improve detection of rule violations, and careful cognitive science to make the system work with real people.

\end{document}
